\documentclass[11pt]{article}

\usepackage{amsmath,amssymb}
\usepackage{xurl}
\usepackage[colorlinks=true,linkcolor=blue,citecolor=blue,urlcolor=blue]{hyperref}

% Shared commands for notes in docs/paper-gaps/.

\DeclareUnicodeCharacter{2080}{\ifmmode{}_{0}\else\textsubscript{0}\fi}
\DeclareUnicodeCharacter{2081}{\ifmmode{}_{1}\else\textsubscript{1}\fi}
\DeclareUnicodeCharacter{2082}{\ifmmode{}_{2}\else\textsubscript{2}\fi}
\DeclareUnicodeCharacter{2099}{\ifmmode{}_{n}\else\textsubscript{n}\fi}

\newcommand{\Lean}{\textsc{Lean}}
\ExplSyntaxOn
\NewDocumentCommand{\leanid}{m}{
  \ifmmode
    \textsf{\detokenize{#1}}
  \else
    \group_begin:
    \urlstyle{sf}
    \tl_set:Nn \l_tmpa_tl {#1}
    \tl_replace_all:Nnn \l_tmpa_tl {\_} {_}
    \exp_args:NV \path \l_tmpa_tl
    \group_end:
  \fi
}
\ExplSyntaxOff
\providecommand{\tr}{\operatorname{tr}}
\newcommand{\tnleanIssueBase}{https://github.com/LionSR/TNLean/issues/}
\newcommand{\tnleanPrBase}{https://github.com/LionSR/TNLean/pull/}
\newcommand{\tnleanDocBase}{https://sirui-lu.com/TNLean/docs/}
\newcommand{\ghissue}[1]{\href{\tnleanIssueBase#1}{GitHub issue \##1}}
\newcommand{\ghpr}[1]{\href{\tnleanPrBase#1}{PR \##1}}
\newcommand{\doclink}[1]{\href{\tnleanDocBase#1.html}{\path{#1}}}

% Dirac notation and the identity operator, as in blueprint/src/macros/common.tex.
% \providecommand keeps note-local definitions authoritative.
\usepackage{dsfont}
\providecommand{\ket}[1]{|#1\rangle}
\providecommand{\bra}[1]{\langle #1|}
\providecommand{\braket}[2]{\langle #1 | #2 \rangle}
\providecommand{\ketbra}[2]{|#1\rangle\!\langle #2|}
\providecommand{\Id}{\mathds{1}}
\providecommand{\one}{\mathds{1}}
\providecommand{\C}{\mathbb{C}}
\providecommand{\MN}[1]{M_{#1}(\C)}
\providecommand{\MD}{M_D(\C)}

% External referencing for paper sources.  The build helper writes sanitized
% auxiliary files under build/paper-aux/ and excludes duplicated source labels.
\usepackage{xr}

% Import a generated paper auxiliary file with a caller-supplied label prefix.
% The candidate paths support builds from docs/paper-gaps/, the repository root,
% and build/paper-aux/ itself.
\newcommand{\importpaperaux}[2]{%
  \IfFileExists{../../build/paper-aux/#2.aux}{%
    \externaldocument[#1]{../../build/paper-aux/#2}%
  }{%
    \IfFileExists{build/paper-aux/#2.aux}{%
      \externaldocument[#1]{build/paper-aux/#2}%
    }{%
      \IfFileExists{#2.aux}{%
        \externaldocument[#1]{#2}%
      }{}%
    }%
  }%
}

% General external reference macros
\newcommand{\paperref}[2]{\ref{#1-#2}}
\newcommand{\papereqref}[2]{(\ref{#1-#2})}

% CPSV16 source (1606.00608 / MPDO-22-12-17-2.tex) external document and shortcuts
\importpaperaux{cpsv16-}{1606.00608/MPDO-22-12-17-2-tnlean}
\newcommand{\cpsvref}[1]{\paperref{cpsv16}{#1}}
\newcommand{\cpsveqref}[1]{\papereqref{cpsv16}{#1}}

% Verdict marker: \gapnote{<kind>}{<status>}. Kind is the policy
% classification (clarification, local-correction, scope-restriction,
% unfaithful, false-source, open-gap); status is open, wip (elimination
% underway), resolved, or historical. Machine-read by the site index;
% prints a verdict line.
\newcommand{\gapnote}[2]{\par\noindent\textbf{Verdict:} #1 (#2).\par}


\title{Homogeneous-Chain Scope for the GNVW Support Algebras,
Adjacent Generation, and Index}
\author{}
\date{2026-08-24}

\begin{document}
\maketitle
\gapnote{scope-restriction}{open}

\section{The site-dependent GNVW construction}

Gross--Nesme--Vogts--Werner consider a chain whose cell algebra at
$y\in\mathbb Z$ is
\begin{align}
    \mathcal A_y&=M_{d(y)}(\mathbb C),
\end{align}
where the positive matrix size $d(y)$ may depend on $y$
\cite[Section~7.2]{Gross2012IndexTheory}.  For a nearest-neighbour
automorphism $\alpha$, set
\begin{align}
    E_x&=\{2x,2x+1\},&
    L_x&=\{2x-1,2x\},&
    P_x&=\{2x+1,2x+2\}.
\end{align}
The locality inclusion is
\begin{align}
    \alpha(\mathcal A_{E_x})
    &\subseteq
      \mathcal A_{L_x}\otimes\mathcal A_{P_x}.
    \label{eq:site_scope_locality}
\end{align}
In the factor order displayed in (\ref{eq:site_scope_locality}), equation
\texttt{RR2x} defines
\begin{align}
    \mathcal R_{2x}
      &=\operatorname{Spp}_{\mathrm L}
        \bigl(\alpha(\mathcal A_{E_x})\bigr)
        \subseteq\mathcal A_{L_x},\\
    \mathcal R_{2x+1}
      &=\operatorname{Spp}_{\mathrm R}
        \bigl(\alpha(\mathcal A_{E_x})\bigr)
        \subseteq\mathcal A_{P_x}.
    \label{eq:site_scope_supports}
\end{align}
These formulas occur at lines 1251--1266 of the arXiv v2 source.  The
site-dependent cell structure and causal-automorphism hypotheses occur at
lines 557--600.  Translation covariance is not used in
(\ref{eq:site_scope_locality}) or (\ref{eq:site_scope_supports}).

\section{Adjacent generation and dimensions}

Let $r(y)>0$ be the matrix size for which
$\mathcal R_y\cong M_{r(y)}(\mathbb C)$.  The adjacent-generation argument
gives the two ordered product identities
\begin{align}
    d(2x)d(2x+1)
      &=r(2x)r(2x+1),                                      \label{eq:site_scope_even}\\
    r(2x+1)r(2x+2)
      &=d(2x+1)d(2x+2).                                    \label{eq:site_scope_odd}
\end{align}
The first is the matrix-size consequence of equation \texttt{AARReven} at
lines 1276--1287; the second is the consequence of equation
\texttt{AARRodd} at lines 1306--1315
\cite{Gross2012IndexTheory}.  Consequently,
\begin{align}
    \frac{r(2x)}{d(2x)}
      &=\frac{d(2x+1)}{r(2x+1)}
       =\frac{r(2x+2)}{d(2x+2)}.
    \label{eq:site_scope_ratio}
\end{align}
Equation (\ref{eq:site_scope_ratio}) is the site-independent multiplicative
index relation.  It is labelled \texttt{dimtransfer} at lines 1317--1329.

\section{The homogeneous specialization}

The present spin-chain algebra fixes one positive on-site matrix size $d$.
Thus
\begin{align}
    d(y)&=d
\end{align}
for every $y$, and (\ref{eq:site_scope_even})--(\ref{eq:site_scope_ratio})
specialize to
\begin{align}
    d^2&=r(2x)r(2x+1),\\
    r(2x+1)r(2x+2)&=d^2,\\
    \frac{r(2x)}{d}
      &=\frac{d}{r(2x+1)}
       =\frac{r(2x+2)}{d}.
    \label{eq:site_scope_homogeneous}
\end{align}
This is the setting used by Cirac--P\'erez-Garc\'ia--Schuch--Verstraete:
their Appendix fixes $\mathcal A_y\cong M_d(\mathbb C)$ and forms the even
support algebra in lines 2320--2329 \cite{Cirac2017MPU}.

The statements presently proved for the fixed-$d$ spin chain establish only
the homogeneous versions of (\ref{eq:site_scope_locality}) and
(\ref{eq:site_scope_supports}), together with their two-sided support
containment.\footnote{Formal declarations:
\leanid{SpinChain.PropagatesWithin.evenPairLocalImage},
\leanid{SpinChain.PropagatesWithin.evenSupportAlgebra},
\leanid{SpinChain.PropagatesWithin.oddSupportAlgebra}, and
\leanid{SpinChain.PropagatesWithin.evenPairLocalImage_le_support_kroneckerSubmodule}
in \doclink{TNLean/QCA/SiteIndexedSupportAlgebra}.}
The adjacent-generation equalities, matrix-factor presentations,
identities in (\ref{eq:site_scope_homogeneous}), and index are not part of
these statements; any fixed-$d$ formulation of those later results remains
subject to the same homogeneous-chain restriction.

This restriction is distinct from the full-matrix-factor restriction in
\path{docs/paper-gaps/gnvw12_support_algebra_full_matrix_scope.tex}.  That note
compares full complex matrix factors with arbitrary finite-dimensional
$C^*$-algebras.  The present note compares a fixed matrix size along the chain
with the site-dependent full matrix sizes of GNVW.

\section{Elimination plan}

Removing the restriction requires a site-dependent algebraic local observable
system with cell algebras $M_{d(y)}(\mathbb C)$, canonical inclusions for
arbitrary finite regions, and its norm completion.  The finite-region
restriction of a causal automorphism must then be written in the ordered
$L_x\times P_x$ coordinates of (\ref{eq:site_scope_locality}).  The same
coefficient definitions give (\ref{eq:site_scope_supports}); the support
commutation, scalar-center, and adjacent-generation arguments then yield
(\ref{eq:site_scope_even}) and (\ref{eq:site_scope_odd}) without imposing
constancy of $d(y)$.  Finally, (\ref{eq:site_scope_ratio}) defines the
site-independent GNVW index.  The fixed-$d$ theorems should thereafter be
derived by specializing the dimension function to the constant value $d$.

Until such a site-dependent observable algebra is available, statements
corresponding to these steps are necessarily restricted to the homogeneous
chain.

\bibliographystyle{alpha}
\bibliography{references}

\end{document}
